\documentclass[11pt]{article}

\usepackage[margin=1in]{geometry}
\usepackage[T1]{fontenc}
\usepackage{lmodern}
\usepackage{amsmath}
\usepackage{amssymb}
\usepackage{booktabs}
\usepackage{hyperref}
\usepackage{xcolor}

\hypersetup{
  colorlinks=true,
  linkcolor=blue!50!black,
  citecolor=blue!50!black,
  urlcolor=blue!50!black,
  pdftitle={Batch Deduplication and Insertion-Time Prevention with a Shared Vector Canopy Index},
  pdfauthor={Denis Robert},
}

\title{Batch Deduplication and Insertion-Time Prevention\\with a Shared Vector Canopy Index}
\author{Denis Robert\\\small \textit{Independent Researcher}}
\date{}

\begin{document}
\maketitle

\begin{abstract}
Entity resolution is commonly deployed in two different workloads: an \emph{in-place}
batch operation that deduplicates an already-populated database, and an \emph{insertion}
operation that stops new duplicates before they land. These are usually treated as separate
problems with separate machinery. This paper argues they should share a single vector index,
and shows how. Dense row embeddings and a vector index serve, first, as the
\emph{cheap phase} of canopy-style blocking for batch deduplication: k-means (or any
vector-native clustering) forms overlapping canopies, and Fellegi--Sunter (Splink) is the
\emph{expensive phase} that scores the candidate pairs within each canopy and joins them into
connected components. The same index and the same canonical record mapping are then reused at
insertion time: an incoming record is embedded, blocked to its canopies, and resolved against
the canonical (deduplicated) store so that an existing record is rejected and a new record is
admitted. We prototype this two-phase batch deduplication on a synthetic Canadian person
database that genuinely contains duplicates, measure cluster-level precision/recall/F1 under
varied canopy overlap, and demonstrate that the canonical index correctly rejects a re-insert
and accepts a new record. We find that embedding-recoverability --- whether a true duplicate
is placed near its mate by the embedding model --- is the binding constraint, and that
canopy overlap is the primary recall lever. We close with the calibration caveat from the
companion Calibration Paradox work: candidate pairs drawn from resemblance-selected
neighbourhoods are even more biased than value-equality collisions, so any prior fitted over
them must be re-anchored and the threshold re-tuned jointly.
\end{abstract}

\section{Introduction}
Entity resolution determines whether two records refer to the same real-world entity. In a
persistent database, two distinct workloads demand resolution:

\begin{enumerate}
  \item \textbf{Batch (in-place).} The database already exists and contains duplicates ---
        accumulated over time, imported from multiple sources, or merged from acquisitions.
        The task is to find all duplicate groups and collapse them, producing a canonical,
        deduplicated database.
  \item \textbf{Insertion.} New records arrive continuously. The task is to prevent
        future duplicates before they are stored --- resolve each incoming record against
        the current database and reject or merge it if it matches.
\end{enumerate}

These are conventionally addressed as separate problems: batch dedup as an offline
blocking-and-linkage job, insertion as a query-time retrieval-and-resolve step. The central
claim of this paper is that \emph{both workloads naturally share one artifact}: a dense
vector index over row embeddings of the database. The index provides the cheap candidate
generation for batch dedup, and, once the database is canonicalised, it is the same index
that insertion-time blocking queries. This yields three concrete benefits:

\begin{itemize}
  \item \textbf{One embedding cost.} Each record is embedded once; neither workload re-embeds
        the population. The repository's persistence design already reloads an index without
        re-embedding~\cite{pipeline}.
  \item \textbf{One structural primitive.} The ``canopy'' that the cheap phase produces for
        batch dedup is the same shape as the top-$k$ neighbourhood an insertion query
        retrieves: a set of resemblance-selected candidates that the expensive Fellegi--Sunter
        stage then scores.
  \item \textbf{One canonical subject.} After batch dedup, the index is rebuilt to contain
        canonical cluster representatives; insertion-time blocking then queries a
        duplicate-free store, so the duplicates that batch dedup removed cannot be re-created.
\end{itemize}

This paper is organised as follows. Section~\ref{sec:use-cases} distinguishes the three
principal entity-resolution workflows --- deduplication, cross-database linking, and merging
--- and locates the shared-index idea within them. Section~\ref{sec:fs} recalls the two-stage
Fellegi--Sunter pipeline from the companion whitepaper~\cite{pipeline}. Section~\ref{sec:canopy}
describes the canopy-style cheap phase --- vector-native clustering with overlapping
multi-assignment --- and the expensive connected-components phase. Section~\ref{sec:dual}
states the dual-use motivation that unifies batch dedup and insertion-time prevention on one
index. Section~\ref{sec:emp} reports a prototype experiment on a synthetic duplicate-bearing
database. Section~\ref{sec:caveats} discusses the embedding-recoverability ceiling and the
calibration caveat. Section~\ref{sec:lit} relates the approach to the literature,
Section~\ref{sec:concl} concludes.

\section{Use cases: deduplication, linking, and merging}
\label{sec:use-cases}
Entity resolution is not a single task. It appears in at least three distinct workflows, and
the differences among them matter for how a vector index is built, queried, and canonicalised.

\begin{enumerate}
  \item \textbf{Deduplication (one database).} A single, already-populated database contains
        duplicate groups representing the same entity. The task is to find every such group
        and collapse it into one canonical record. This is the in-place batch workload of
        Sections~\ref{sec:canopy}--\ref{sec:emp}: a single input population is both the
        candidate source and the canonical subject, and the output is a deduplicated database.
  \item \textbf{Linking (two or more databases).} Two (or more) separately maintained
        databases hold records that may refer to the same entities, and the task is to
        establish the \emph{correspondence} between them --- ``this record in database A is
        the same person as that record in database B'' --- \emph{without} collapsing either
        database. The matched pairs are used to join analytics, to reconcile identifiers, or
        to propagate updates between systems, while each source keeps its own records. This
        is the \texttt{link} (or \texttt{link\_only}) variant of the Fellegi--Sunter settings:
        two input datasets with distinct ``source dataset'' labels are compared against each
        other, and the output is a set of match edges, not collapsed clusters.
  \item \textbf{Merging (collapse across databases).} The linking step is followed by a
        consolidation step: two databases are linked, then their matched records are merged
        into a single combined entity set, with conflicting fields resolved and duplicates
        eliminated. Merging therefore \emph{reduces to deduplication}: once the cross-database
        correspondence is known (or once the sources are physically concatenated), the problem
        is again to deduplicate a single combined population and produce canonical records.
\end{enumerate}

These three are related but not interchangeable. \emph{Deduplication} and \emph{merging} both
end in canonical records, so they share the connected-components machinery; the only
difference is whether the input is one population (dedup) or several concatenated ones
(merge). \emph{Linking} is genuinely distinct: the output is a set of cross-database edges
that \emph{preserves} both source databases rather than collapsing them, so there is no single
canonical index to maintain and no cluster to consolidate --- only a persistent mapping between
record identifiers across sources.

The shared-index thesis of this paper extends to linking with one adaptation. The cheap phase
--- vector clustering into overlapping resemblance canopies --- is index-agnostic to whether
the canopies span one database or several: the candidate pairs are simply those records (from
any source) that share a canopy, and the expensive phase scores each pair under the same
Fellegi--Sunter comparisons. For \emph{linking}, the vector index and its canopies can be
built over the concatenation of all sources (one shared embedding space), and the
\emph{link\_only} settings --- which attach a ``source dataset'' column and emit pairwise
edges instead of connected-component clusters --- are used in place of \texttt{dedupe\_only}.
For \emph{merging}, the same link edges feed a consolidation pass that produces the combined
deduplicated index. The dual-use property of Section~\ref{sec:dual} (one canonical index
serving batch and insertion workloads) is specific to deduplication and merging, where a
canonical subject exists; linking, by contrast, maintains a \emph{correspondence}, and its
``insertion-time'' analogue is matching a newly arriving record in one source against the
other source's index while still preserving both.

\section{Two-stage Fellegi--Sunter pipeline}
\label{sec:fs}
The companion whitepaper~\cite{pipeline} describes a two-stage entity-resolution pipeline:
dense row embeddings (sentence-transformers/all-MiniLM-L6-v2) persisted in a FAISS
\emph{IndexFlatIP} store, with each query embedded and blocked to its top-$k$ cosine
neighbours, after which Splink computes a match probability over explicit field comparisons
(Jaro--Winkler for names and address, a date-of-birth comparison, an email comparison) and
returns a match when the probability exceeds a threshold $\tau$. The pipeline is designed for
\emph{insertion-time} resolution: given one query record, retrieve a small candidate set and
decide. The present paper generalises the same two stages to the batch workload.

\section{Canopy-style batch deduplication}
\label{sec:canopy}

\subsection{Cheap phase: vector-native clustering as canopies}
For batch deduplication the object is not a single query but the whole database. The cheap
phase groups the stored vectors into candidate sets --- \emph{canopies}. The canonical recipe
is canopy clustering~\cite{mccallum2000}: a cheap distance measure places each record into
overlapping candidates (a wide capture radius gathers members, a tighter criterion stops
growth), and an expensive similarity is applied only within the resulting canopies. The
overlap is deliberate: it prevents a true duplicate from being missed purely because it
straddles a cluster boundary.

A vector index supplies a natural cheap phase. In the prototype we use FAISS k-means as the
coarse quantizer and \emph{multi-assignment} for overlap: each record is assigned to its top
$m$ centroids rather than a single one. With $m>1$ the canopies overlap, and the candidate
pairs are exactly the record pairs that share a centroid. This is a direct vector analogue of
the two-threshold canopy procedure: the number of centroids $C$ and the assignment depth $m$
stand in for the capture and stop thresholds. Because vectors are L2-normalized and k-means
operates on cosine-consistent geometry, the canopies reflect embedding similarity.

\subsection{Expensive phase: Fellegi--Sunter over canopies + connected components}
Each canopy is a shared block. Fellegi--Sunter (Splink) scores every candidate pair within a
canopy under the field comparisons and a decision threshold $\tau$. Pairs scoring at or above
$\tau$ become match edges, and connected components over those edges yield the duplicate groups
(``clusters''). This is batch deduplication as transitive closure: if record $A$ matches $B$
and $B$ matches $C$, all three are placed in one canonical cluster even if $A$ and $C$ were
never directly compared. Splink's clustering utilities compute exactly this on the predicted
edges.

Implementation note: Splink's \texttt{dedupe\_only} linker requires a single input table, so
rather than passing a separate edge table we encode canopy membership as wide
\texttt{anchor\_$j$} columns (one per centroid assignment) on the node table and give Splink
one blocking rule per column. This reconstructs exactly the canopy candidate pairs while
satisfying the single-table constraint.

\subsection{The candidate-pair cost surface}
The cheap phase controls the cost of the expensive phase. With $C$ canopies and assignment
depth $m$, a record belongs to up to $m$ canopies, and each canopy of size $s$ yields
$s(s-1)/2$ candidate pairs. The total candidate-pair budget therefore grows with cluster
size (``quadratically in the number of records per canopy'') and with $m$ (through overlap).
Raising $C$ shrinks canopies and lowers the budget but risks missing cross-boundary
duplicates; raising $m$ recovers recall that hard partitions lose but inflates the budget.
This precision--recall--cost trade-off of $(C,m)$ is the batch analogue of the online $k$
trade-off in the whitepaper~\cite{pipeline}.

\section{Dual use: batch dedup and insertion-time prevention share one index}
\label{sec:dual}
The pragmatic motivation for using a vector index as the cheap phase is that the same index
serves both workloads:

\emph{One embedded index, two workloads.} The batch pass runs offline over the populated
index: cluster (cheap phase), run Fellegi--Sunter within canopies, collapse connected
components into canonical records, and rebuild the index to hold one canonical representative
per cluster. The insertion pipeline then queries that \emph{same} canonical index: an
incoming record is embedded, blocked to its top-$k$ canopies, and resolved against canonical
cluster representatives. The embedding cost is paid once; the persistence design in the
repository~\cite{pipeline} (\texttt{index.faiss} + metadata, reloaded without re-embedding)
makes this reuse operationally natural.

Two consistency requirements make the claim non-trivial rather than automatic:

\begin{enumerate}
  \item \textbf{The index must reflect the deduplicated database.} Raw batch dedup over the
        original index sees duplicated vectors as separate records. After collapsing
        duplicates, the index must be rebuilt from the canonical clusters (deduplicated
        vectors, remapped metadata) so insertion-time blocking queries a duplicate-free
        store --- otherwise the index still ``contains'' the very duplicates it was meant to
        prevent. The store's \texttt{add}/\texttt{delete}/\texttt{reindex} semantics are the
        seam for this.
  \item \textbf{Canonical representatives.} For insertion, the resolvable subject is the
        consolidated cluster (e.g.\ one representative per cluster), not every member of the
        original duplicate set. The batch pass must emit a vector-position $\to$ canonical-cluster
        mapping, which is exactly what insertion-time blocking consumes.
\end{enumerate}

\section{Prototype experiment}
\label{sec:emp}
We implement \texttt{scripts/experiment\_batch\_dedup.py} in the companion repository, which
prototypes the full loop (cheap clustering $\to$ Splink dedupe/connected components $\to$
canonicalisation $\to$ insertion check) on a synthetic Canadian person database that
genuinely contains duplicate pairs.

\subsection{Setup}
A base population of $N$ synthetic Canadian person records (Faker with a custom Canadian
address provider, 30\% independent address/email missingness) has a near-duplicate ``twin''
inserted for a $3\%$ fraction, so the reference index genuinely contains matching records
(the same construction as the duplicate-bearing benchmark of the whitepaper~\cite{pipeline}).
Two twin styles are used:

\begin{description}
  \item[identical] the twin is an exact copy of the base record --- a positive control that
        trivially exercises the full canopy-dedupe machinery; and
  \item[varied] the twin is a perturbed near-duplicate (typos, address/email variation),
        exposing whether the embedding itself recovers the pair.
\end{description}

The cheap phase runs FAISS k-means with $C$ centroids and assignment depth $m$; Splink scores
the within-canopy pairs at $\tau=0.85$ (untrained default $m/u$, prior
$10^{-4}$); connected components yield cluster ids. Cluster-level precision/recall/F1 are
measured against the known ground-truth duplicate pairs. In a validation environment
(a 3{,}000-record base), the identical-twin control at $C=128, m=2$ reached recall $1.0$ and
F1 $0.968$ (precision $0.938$), with 90/90 twins recovered by the cheap phase.

\subsection{Overlap is the recall lever}
Table~\ref{tab:cheap-recall} reports cheap-phase (canopy) twin recall against assignment
depth $m$. Raising $m$ recovers materially more true duplicates: at $m=2$ only a fraction of
varied twins share a canopy, whereas at $m=8$ all 90 are captured. This is the direct
prediction of Section~\ref{sec:canopy}: hard partitions lose cross-boundary duplicates, and
overlap recovers them. It also confirms the central empirical claim of this paper ---
\emph{embedding-recoverability is the binding constraint}: with the MiniLM embedding, exact
twins are recovered by the embedding and the pipeline; varied twins are recovered only once
aggressive overlap compensates.

\begin{table}[h]
\centering
\caption{Cheap-phase (canopy) recall of the 90 base/twin duplicate pairs as a function of
assignment depth $m$ ($C=128$ centroids).}
\label{tab:cheap-recall}
\begin{tabular}{lcc}
\toprule
$m$ & identical twins recovered & varied twins recovered \\
\midrule
2 & 90 / 90 & 6 / 90 \\
8 & 90 / 90 & 90 / 90 \\
\bottomrule
\end{tabular}
\end{table}

The varied-twin result at $m=8$ also carries forward to end-to-end linkage: cluster metrics
were F1 $0.90$ (precision $0.90$, recall $0.90$) at $\tau=0.85$, with 81 true positives and 9
false negatives among the 90 ground-truth duplicate pairs. The false negatives are exactly
the cheap-phase misses that overlap at $m=8$ still cannot recover for the noisy records.
These are synthetic-sample measurements, not a general guarantee; they serve to demonstrate
the mechanism and its data-dependent sensitivity, consistent with the tone of the companion
work~\cite{pipeline,paradox}.

\subsection{Insertion-time reuse}
After canonicalisation (one representative per cluster, twins removed, index rebuilt), the
prototype resolves two incoming records against the canonical store: an identical re-insert of
a base record, and a genuinely unrelated new record. The re-insert is retrieved and matched at
match probability $1.0$ against the canonical position of the original record and rejected as a
duplicate; the unrelated record retrieves candidates but no match above threshold and is
accepted as new. This demonstrates the shared-index loop: the same canonical index that
batch-dedup built also serves insertion-time duplicate prevention.

\section{Discussion and caveats}
\label{sec:caveats}

\subsection{Embedding-recoverability is the ceiling}
The dominant failure mode is not the linkage stage but the cheap phase: if the embedding does
not place a true duplicate near its mate, no canopy overlap or threshold tuning can recover it.
This mirrors the whitepaper's ``blocking recall is the binding constraint'' finding for the
online pipeline~\cite{pipeline}. In this prototype the MiniLM embedding recovers exact twins
reliably but is sensitive to address edits for varied twins. The implication is that
improvements to the cheap phase (richer embeddings, multi-view retrieval, or a dedicated
canopy-aware index) are the lever that raises the batch-dedup ceiling, not more aggressive
Splink settings.

How much this sensitivity matters depends on the data pipeline upstream of entity resolution.
The operational problem that originally motivated this line of work supplies person records
whose addresses are already standardized against an external postal reference engine before
they reach either the embedding stage or the linkage stage, so a variant record's address is
already in a canonical street/unit/city/postcode form rather than an arbitrary free-text form.
Under that precondition the address edit that dominated the embedding drift in this synthetic
measurement --- where variation is injected into the raw address string --- is largely
removed, so the address-driven recoverability loss is not expected to recur in that setting.
The findings here are therefore best read as a lower bound on what a \emph{generic}
embedding sees when fed un-standardised address text. In the general case --- where addresses
arrive un-parsed and inconsistent across sources --- the measured sensitivity indicates that
a better embedding strategy could pay off: either richer or domain-adapted encoders that
weigh the more stable identity fields (name, birth date) ahead of the volatile address field,
or a multi-view retrieval that embeds identity, contact, and address fields separately so that
a changed address does not pull a true duplicate out of the canopy. The concrete choice remains
data-dependent and should be evaluated with the same cheap-phase-recall protocol used here,
on the actual address distribution of the deployment.

\subsection{Address fluidity over time}
A further, distinct reason the address is a fragile blocking signal is that it is \emph{time
limited}: unlike a date of birth, which is invariant, the address changes with the entity's
movement, and it does so in an uncontrolled way. The informativeness of the address therefore
depends on the age gap between the two records, not only on how it was written. Within a short
window (two records created close together at the same address, e.g.\ under the typical
residency period of roughly a decade) the address is stable and highly informative. Once
records are separated by more than that window --- the same entity who has since \emph{moved} ---
the address no longer agrees, and blocking on it becomes actively harmful: the address equality
that would bind two stable-at-a-standstill records is absent, and a true duplicate with a
changed address may fall out of a small $k$. Email is more stable than the address across time,
and names change only in controlled ways (nicknames, abbreviations, marriage), so date of birth
and normalized name remain the reliable invariant keys across long gaps.

We measure this with a dedicated script, \texttt{experiment\_temporal\_gap.py}, which creates
duplicate pairs at a controlled age gap: for a ``short'' gap the duplicate keeps the same
address (stable), and for a ``long'' gap (beyond the residency window) the duplicate has
integrally \emph{moved} to a different address while retaining name, date of birth, and email.
Table~\ref{tab:temporal-recall} reports top-$k$ recall (MiniLM, 6{,}000-record base, 180
pairs). In the short gap the ``full`` and ``identity`` views already reach $1.0$ at $k{=}1$;
in the long gap the \emph{contact} (address + email) view collapses to $0.22$ while the
\emph{identity} view (first, last, date of birth) stays at $1.0$.

The same failure mode is confirmed on \emph{real} person data. We used NC voter registration
snapshots of a single county (Wake) across 2012--2026, joined on the stable \emph{voter id},
giving the same voter at multiple points in time with real addresses and real capture dates.
Two complementary studies were run. First, as a \emph{mechanism probe}, a single pair of
snapshots fourteen years apart split by whether the voter's street changed shows that a changed
address destroys retrieval: on a balanced 10{,}000-pair sample (MiniLM, recall at $k{=}1$) the
address-only ``contact'' view recovers $0.97$ of stayed voters but only $0.03$ of moved voters,
the single-vector ``full'' view drops from $0.998$ (stayed) to $0.83$ (moved), and the
invariant ``identity'' view stays at $0.92$ among movers. Second, as the \emph{honest, primary}
methodology, we used \emph{staggered arrival}: for each voter we picked a random historical
snapshot as their entry row and matched their 2026 row against it, so the reference rows span a
genuine varied age-gap distribution (two to fourteen years) with no \emph{moved} label used in
retrieval or selection. On 3{,}000 voters (recall at $k{=}1$): ``full'' $0.985$, ``identity''
$0.991$, ``contact'' $0.872$, ``gap-weighted'' $0.967$, and a hard ``two-tier'' bucket $0.969$.
Because these gaps (max 14 yr) all fall \emph{within} the county's fitted residency window
(below), the address remains mostly informative and ``contact'' is degraded but not collapsed;
``identity'' is the most robust single view, and the two age-aware methods are near-tied. The
decay's distinctive payoff sits in the beyond-residency regime, which a single county has too few
pairs to exercise. The \emph{residency} was estimated by fitting a censored Weibull to
registration durations in the same snapshots ($\texttt{estimate\_residency.py}$): shape
$k\approx1.18$ and mean $\bar\tau\approx20.6$ yr (by both MLE and KM-regression), with a median
of roughly 16 years. Since $k\approx1$, the exponential form $e^{-\Delta t/T}$ with
$T=\bar\tau\approx20.6$ yr is a close model, replacing the synthetic assumption $T{=}10$.

The \emph{gap-weighted} view shows the recommended resolution. Instead of a hard
recent/beyond-residency cutoff, it ages the address \emph{continuously}: the identity and
address vectors are each embedded and searched, and the candidate re-ranking is
\begin{equation}
\text{combined} = s_{\text{identity}} + w(\Delta t)\,s_{\text{address}},\qquad
w(\Delta t)=e^{-\Delta t/T},
\label{eq:gapweight}
\end{equation}
where $\Delta t$ is the age gap between the two records and $T$ the residency timescale, now
estimated from the data ($\approx20.6$ yr). The identity score always counts at full weight; the
address score is down-weighted exponentially as records age apart, so an old address cannot drag a
true long-gap duplicate below top-$k$ while it still boosts recent pairs. This is a smooth
generalisation of the binary rule, and the address weight is applied at \emph{score-fusion time}
(after retrieval, where the pair-wise gap is known) rather than baked into a single precomputed
vector. In the synthetic long-gap case (Table~\ref{tab:temporal-recall}) the gap-weighted view
recovers to $1.0$ by $k{=}5$ while still exploiting the address for recent records; on the
staggered real data (above) it lands between ``full'' and ``identity'' at $k{=}1$, consistent
with a mostly-informative address within the residency window. The claimed benefit is
specifically for pairs spanning beyond $\sim T$ years, where a preserved address would otherwise
mislead.

\begin{table}[h]
\centering
\caption{Top-$k$ blocking recall of true duplicates by record age gap (MiniLM; $T{=}10$yr).}
\label{tab:temporal-recall}
\begin{tabular}{lccc}
\toprule
view & short-gap $k{=}1$ & long-gap $k{=}1$ & long-gap $k{=}5$ \\
\midrule
full         & 1.00 & 0.95 & 0.99 \\
identity     & 1.00 & 1.00 & 1.00 \\
contact      & 0.89 & 0.22 & 0.23 \\
gap-weighted & 0.91 & 0.99 & 1.00 \\
\bottomrule
\end{tabular}
\end{table}

Two implications follow. First, the address is made a \emph{decaying}, not a constant,
retrieval signal: weighting it by $e^{-\Delta t/T}$ in a multi-index fusion preserves its
value for recent records while guarding against the long-gap failure mode that a constant
address weight introduces. Second, the same age-gap decay should apply in the Fellegi--Sunter
stage --- the address comparison weighted strongly within the residency window and
exponentially down-weighted beyond it --- rather than a single address comparison applied to
all pairs. This mirrors the whitepaper's finding that address is informative, with the
important timescale caveat that it is informative \emph{conditionally}, on records being
temporally near each other; the decay rate $T$ must be set from the deployment's observed
residency distribution rather than assumed. On the co-temporal synthetic population of
Section~\ref{sec:emp} (all records sharing one generation date) the temporal caveat is
invisible, so the measured address sensitivity there should be read as the same-window case,
not the general one.

\subsection{Field-position sensitivity of the embedding}
A related question is whether the \emph{ordering} of fields inside the serialized row affects
the embedding's canopy behaviour --- that is, whether fields earlier in the string carry more
weight than later ones. We probed this with a dedicated script,
\texttt{experiment\_field\_ordering.py}, using three measurements.

First, a \emph{positional gradient}: a fixed-content probe sentence plus a single unique
marker token moved to each position was embedded and compared with a no-marker baseline. The
cosine distance of the marker is strongly front-loaded --- $0.083$ at position~0 versus roughly
$0.012$ mid-string --- then plateaus, with only a small rise at the final position. Earlier
fields are therefore more weighty, but the effect is confined to the first couple of tokens
rather than increasing monotonically across the whole string. This is consistent with the
mean-pooled, positionally-embedded architecture of all-MiniLM-L6-v2.

Second, a \emph{permutation-invariance} index: for records with all five fields present, the
mean pairwise cosine across ~24 field orderings of the same record is $0.973$. Reordering thus
moves the embedding by roughly $2.7\%$ --- a real but modest perturbation, smaller than the
token-level gradient, because whole labeled fields saturate attention less than isolated
tokens.

Third, the \emph{canopy impact}: re-embedding the reference population under several orderings
and clustering each with the same canopy parameters ($C=128$, $m=2$) changes cheap-phase recall
of the varied twins in a measurable but bounded way. Leading with the long, discriminative
\emph{address} field recovered the most twins ($0.978$, 88/90) at the \emph{lowest} candidate
budget; the default order recovered $0.944$ (85/90); and \emph{identity-first} recovered the
fewest ($0.922$, 83/90). Field ordering is therefore a small, mostly free lever on cheap-phase
recall --- but this result is data-dependent, here favouring the field that is most informative
on this population, and it sits within the embedding-recoverability ceiling discussed above.
As with any embedding choice, the winning order should be selected on deployment data with the
same cheap-phase-recall protocol, not assumed generally.

\subsection{Calibration caveat from the companion work}
The companion Calibration Paradox paper~\cite{paradox} shows that refitting Fellegi--Sunter
parameters on resemblance-selected candidate pairs inflates the fitted prior: candidate sets
drawn from records that already resemble each other concentrate the non-match mass around
near-collisions, so a free $\hat\lambda$ drifts upward and the classes overlap. Canopy-based
batch dedup makes this \emph{worse}, not better: embedding-neighbourhood candidate pairs are
even more resemblance-biased than value-equality collisions, so a prior fitted over them is
expected to drift beyond the $7.2\times10^{-3}$ observed with value blocking. The practical
rule from that paper applies a fortiori here: \emph{anchor the prior} (default or
blocking-adjusted), fit $m/u$ with the prior pinned, and select $\tau$ jointly on a
validation set.

\subsection{Operational considerations}
Batch dedup and insertion-time prevention have different latencies and cadences, and the
shared index must be maintained coherently: deletes and updates from dedup must propagate to
the canonical store before insertion-time blocking relies on it. In practice the operating
loop is ``deduplicate in batches, prevent at the boundary continuously,'' with periodic
re-clustering or incremental index maintenance as the population evolves. This is operational
scaffolding rather than an algorithmic limit, and it is the same seam identified in the
whitepaper's discussion of index maintenance.

\section{Related work}
\label{sec:lit}
Canopy clustering was introduced by McCallum, Nigam, and Ungar for efficient blocking before
expensive reference matching~\cite{mccallum2000}; the cheap-then-expensive structure is
exactly what a vector index plus Fellegi--Sunter reproduces. Embedding-based blocking has been
used to build candidate sets for downstream matchers and training in DeepER~\cite{ebreheem2018},
and similarity-based (including vector) blocking has been compared directly with
exact-match blocking in the survey spirit of Zeakis et al.~\cite{zeakis2023} and
DeepBlocker~\cite{thirumuruganathan2021}. Learned blocking functions that select candidate
generation from labelled seeds (AutoBlock~\cite{zhang2020autoblock}) are data-driven analogues
of canopy selection. The novel emphasis here is the \emph{dual use}: the same embedded,
canonicalised index serves both the batch-dedup cheap phase and insertion-time prevention, a
property the blocking literature treats only implicitly as ``the index feeds the matcher.'' To
the best of our knowledge the explicit two-workload sharing of one canonical vector index is
under-emphasised in prior work.

\section{Conclusion}
\label{sec:concl}
A vector index unifies batch deduplication and insertion-time prevention. For batch dedup, the
index supplies the canopy cheap phase (vector-native clustering with overlapping
multi-assignment), and Fellegi--Sunter is the expensive phase that scores within-canopy
candidate pairs and joins them into connected components. Once the database is canonicalised,
the same index resolves incoming records to stop future duplicates. A prototype on a synthetic
duplicate-bearing database demonstrates the full loop, confirms that overlap is the recall
lever, and shows that the canonical index rejects re-inserts and accepts new records. The two
caveats that govern practical use are the embedding-recoverability ceiling and the
prior-inflation caveat from the companion Calibration Paradox work; together they dictate that
the cheap phase, not the linkage stage, is where batch-dedup quality is won or lost, and that
any fitted prior must be re-anchored and the threshold re-tuned jointly.

\begin{thebibliography}{9}
\bibitem{pipeline}
Denis Robert.
\newblock A two-stage entity resolution pipeline with FAISS blocking and Splink matching.
\newblock Technical report, Independent Researcher, 2026.

\bibitem{paradox}
Denis Robert.
\newblock The calibration paradox in Fellegi--Sunter processes.
\newblock Technical report, Independent Researcher, 2026.

\bibitem{mccallum2000}
Andrew McCallum, Kamal Nigam, and Lyle H. Ungar.
\newblock Efficient clustering of high-dimensional data sets with application to reference matching.
\newblock In Proceedings of the Sixth ACM SIGKDD International Conference on Knowledge Discovery and Data Mining, pages 169--178, 2000.
\newblock \url{https://doi.org/10.1145/347090.347123}.

\bibitem{ebreheem2018}
Muhammad Ebraheem, Saravanan Thirumuruganathan, Shafiq Joty, Mourad Ouzzani, and Nan Tang.
\newblock Distributed representations of tuples for entity resolution.
\newblock Proceedings of the VLDB Endowment, 11(11):1454--1467, 2018.
\newblock \url{https://doi.org/10.14778/3236187.3236198}.

\bibitem{zeakis2023}
Alexandros Zeakis, George Papadakis, Dimitrios Skoutas, and Manolis Koubarakis.
\newblock Pre-trained embeddings for entity resolution: An experimental analysis.
\newblock Proceedings of the VLDB Endowment, 16(9):2225--2238, 2023.
\newblock \url{https://doi.org/10.14778/3598581.3598594}.

\bibitem{thirumuruganathan2021}
Saravanan Thirumuruganathan, Han Li, Nan Tang, Mourad Ouzzani, Yash Govind, Derek Paulsen, Glenn Fung, and AnHai Doan.
\newblock Deep learning for blocking in entity matching.
\newblock Proceedings of the VLDB Endowment, 14(11):2459--2472, 2021.
\newblock \url{https://doi.org/10.14778/3476249.3476294}.

\bibitem{zhang2020autoblock}
Wei Zhang, Hao Wei, Bunyamin Sisman, Xin Luna Dong, Christos Faloutsos, and David Page.
\newblock AutoBlock: A hands-off blocking framework for entity matching.
\newblock In Proceedings of the 13th ACM International Conference on Web Search and Data Mining, pages 744--752, 2020.
\newblock \url{https://doi.org/10.1145/3336191.3371813}.
\end{thebibliography}

\section*{Use of Artificial Intelligence}
During the preparation of this work the author used an AI coding and writing assistant (the
DeepSeek V4 Flash 0731 large language model) to assist with drafting and editing the
manuscript and preparing the experiments and the accompanying prototype code. The tool was
not listed as an author. The design and reported results were reviewed and verified by the
author, who takes full responsibility for the content, including any remaining errors and/or
omissions.

\end{document}
